\documentclass[14pt, a4paper]{extarticle}
\usepackage[utf8]{inputenc}
\usepackage[english,ukrainian]{babel}
\usepackage{amsmath, amssymb, amsfonts, amsthm}
\usepackage{geometry}
\usepackage{graphicx}
\usepackage{hyperref}
\usepackage{enumitem}
\usepackage{array}
\usepackage{booktabs}
\usepackage{xcolor}
\usepackage{colortbl}
\usepackage{longtable}
\usepackage{multirow}
\usepackage{float}

\usepackage{fontspec}
\setmainfont{Times New Roman}
\usepackage[T2A]{fontenc}

\geometry{top=2cm, bottom=2cm, left=2.5cm, right=2cm}

\hypersetup{
    colorlinks=true,
    linkcolor=blue,
    urlcolor=blue,
}

\numberwithin{equation}{section}

\newtheorem{definition}{Визначення}[section]
\newtheorem{theorem}{Теорема}[section]
\newtheorem{proposition}{Твердження}[section]

\title{МОДЕЛІ ДИСТАНЦІЙНОЇ ІДЕНТИФІКАЦІЇ ПАРАМЕТРІВ ДИНАМІЧНИХ ОБ'ЄКТІВ}
\author{Наукове дослідження}
\date{\today}

\begin{document}

\maketitle

\tableofcontents

\newpage

\section{МОДЕЛІ ДИСТАНЦІЙНОЇ ІДЕНТИФІКАЦІЇ ПАРАМЕТРІВ ДИНАМІЧНИХ ОБ'ЄКТІВ}

У цьому розділі розглядаються математичні моделі для дистанційної ідентифікації параметрів динамічних об'єктів. Основна увага приділяється концептуальним засадам побудови моделей, методам ідентифікації параметрів за відеоданими, використанню оптичного потоку для визначення кінематичних характеристик, а також підходам до об'єднання різних моделей і методів.

\subsection{Концепція моделей дистанційної ідентифікації}

\subsubsection{Основні поняття та визначення}

\begin{definition}
\textbf{Дистанційна ідентифікація динамічного об'єкта} — це процес визначення просторових, кінематичних та динамічних характеристик об'єкта за даними відеоспостереження без безпосереднього контакту з об'єктом.
\end{definition}

\begin{definition}
\textbf{Модель дистанційної ідентифікації} — це математичний апарат, що встановлює зв'язок між спостережуваними відеоданими $V(x,y,t)$ та параметрами об'єкта $\Theta = \{\theta_1, \theta_2, \ldots, \theta_n\}$.
\end{definition}

Загальна концептуальна модель дистанційної ідентифікації може бути представлена у вигляді:

\begin{equation}
\boxed{\Theta = \mathcal{F}(V(x,y,t), \mathcal{M}, \mathcal{A})} \tag{3.1}
\end{equation}

де:
\begin{itemize}
\item $\Theta$ — множина параметрів об'єкта, що ідентифікуються;
\item $V(x,y,t)$ — відеодані (послідовність кадрів);
\item $\mathcal{M}$ — множина математичних методів обробки;
\item $\mathcal{A}$ — множина архітектур нейронних мереж.
\end{itemize}

\subsubsection{Ієрархія параметрів динамічних об'єктів}

Параметри динамічних об'єктів можна класифікувати за рівнем абстракції:

\begin{enumerate}
    \item \textbf{Геометричні параметри} — форма, розміри, просторове положення.
    \item \textbf{Кінематичні параметри} — швидкість, напрямок руху, траєкторія, прискорення.
    \item \textbf{Динамічні параметри} — маса, імпульс, енергія (оцінюються опосередковано).
\end{enumerate}

Математично ця ієрархія може бути представлена як:

\begin{equation}
\Theta = \{\Theta_{\text{geom}}, \Theta_{\text{kin}}, \Theta_{\text{dyn}}\} \tag{3.2}
\end{equation}

\begin{align}
\Theta_{\text{geom}} &= \{(x_c, y_c, z), (w, h), \text{shape}\} \tag{3.3} \\
\Theta_{\text{kin}} &= \{v, \theta, T(t), a\} \tag{3.4} \\
\Theta_{\text{dyn}} &= \{m, p, E\} \tag{3.5}
\end{align}

\subsubsection{Архітектура системи ідентифікації}

Система дистанційної ідентифікації складається з наступних модулів:

\begin{equation}
\mathcal{S} = \{\mathcal{D}, \mathcal{F}, \mathcal{Z}, \mathcal{P}, \mathcal{V}\} \tag{3.6}
\end{equation}

де:
\begin{itemize}
\item $\mathcal{D}$ — модуль виявлення об'єктів (DETR/YOLO);
\item $\mathcal{F}$ — модуль оптичного потоку (Farneback/Lucas-Kanade/Horn-Schunck/FlowNet/RAFT);
\item $\mathcal{Z}$ — модуль оцінки глибини (MiDaS/DPT/GeoNet);
\item $\mathcal{P}$ — модуль обчислення параметрів;
\item $\mathcal{V}$ — модуль візуалізації.
\end{itemize}

\begin{figure}[H]
\centering
\begin{tabular}{|c|c|c|c|}
\hline
\textbf{Вхід} & \textbf{Модулі обробки} & \textbf{Вихід} \\
\hline
Відеопотік & $\mathcal{D} \rightarrow \mathcal{F} \rightarrow \mathcal{Z} \rightarrow \mathcal{P} \rightarrow \mathcal{V}$ & Параметри $\Theta$ \\
\hline
\end{tabular}
\caption{Архітектура системи дистанційної ідентифікації}
\label{fig:architecture}
\end{figure}

\subsection{Моделі ідентифікації параметрів за відеоданими}

\subsubsection{Модель виявлення об'єктів на основі DETR}

Модель DETR (Detection Transformer) базується на архітектурі трансформера та забезпечує кінцево-кінцеве виявлення об'єктів. Основна формула механізму уваги:

\begin{equation}
\boxed{\text{Attention}(Q, K, V) = \text{softmax}\left(\frac{QK^T}{\sqrt{d_k}}\right)V} \tag{3.7}
\end{equation}

Багатоголова увага дозволяє моделювати різні типи залежностей:

\begin{equation}
\boxed{\text{MultiHead}(Q, K, V) = \text{Concat}(\text{head}_1, \ldots, \text{head}_h) W^O} \tag{3.8}
\end{equation}

\begin{equation}
\text{head}_i = \text{Attention}(QW_i^Q, KW_i^K, VW_i^V) \tag{3.9}
\end{equation}

Оптимальне парування між передбаченнями та істинними об'єктами знаходиться за допомогою угорського алгоритму:

\begin{equation}
\hat{\sigma} = \arg\min_{\sigma \in \mathfrak{S}_N} \sum_{i=1}^N \mathcal{L}_{\text{match}}(y_i, \hat{y}_{\sigma(i)}) \tag{3.10}
\end{equation}

Функція втрат DETR має вигляд:

\begin{equation}
\boxed{\mathcal{L}_{\text{Hungarian}}(y, \hat{y}) = \sum_{i=1}^N \left[ -\log \hat{p}_{\hat{\sigma}(i)}(c_i) + \mathbf{1}_{\{c_i \neq \varnothing\}} \mathcal{L}_{\text{box}}(b_i, \hat{b}_{\hat{\sigma}(i)}) \right]} \tag{3.11}
\end{equation}

\subsubsection{Модель оцінки глибини MiDaS}

Модель MiDaS призначена для монокулярної оцінки глибини. Вона розв'язує проблему відносної глибини:

\begin{equation}
d_{\text{pred}}(p) = \alpha \cdot d_{\text{true}}(p) + \beta \tag{3.12}
\end{equation}

Багатомасштабна функція втрат забезпечує навчання на різних масштабах:

\begin{equation}
\mathcal{L} = \frac{1}{N} \sum_{i=1}^N \frac{1}{2\sigma_i^2} \mathcal{L}_i + \log \sigma_i \tag{3.13}
\end{equation}

Інваріантна до масштабу втрата дозволяє навчатися на даних з різними масштабами глибини:

\begin{equation}
\boxed{\mathcal{L}_{\text{SI}} = \frac{1}{N} \sum_{i,j} \left[ (\log y_i - \log \hat{y}_i) - (\log y_j - \log \hat{y}_j) \right]^2} \tag{3.14}
\end{equation}

\subsubsection{Модель DPT для щільної оцінки глибини}

DPT (Dense Prediction Transformer) поєднує трансформерну архітектуру з механізмом уваги:

\begin{equation}
z_l = \text{MSA}(\text{LN}(z_{l-1})) + z_{l-1}, \quad l=1,\ldots,L \tag{3.15}
\end{equation}

Вхідні токени формуються як:

\begin{equation}
z_0 = [x_{\text{class}}; x_p^1 E; x_p^2 E; \ldots; x_p^N E] + E_{\text{pos}} \tag{3.16}
\end{equation}

\subsection{Визначення параметрів методами оптичного потоку}

\subsubsection{Модель швидкості об'єкта}

Швидкість об'єкта визначається як середня величина оптичного потоку в області об'єкта $\Omega$:

\begin{equation}
\boxed{v = \frac{1}{N_\Omega} \sum_{(x,y) \in \Omega} \|F(x,y)\|} \tag{3.17}
\end{equation}

де $F(x,y) = [u(x,y), v(x,y)]^T$ — вектор оптичного потоку.

В інтегральній формі:

\begin{equation}
v = \frac{1}{|\Omega|} \iint_{\Omega} \sqrt{u(x,y)^2 + v(x,y)^2} \, dx\,dy \tag{3.18}
\end{equation}

\subsubsection{Модель напрямку руху}

Напрямок руху об'єкта обчислюється через векторне усереднення кутів:

\begin{equation}
\theta(x,y) = \arctan\left( \frac{v(x,y)}{u(x,y)} \right) \tag{3.19}
\end{equation}

\begin{equation}
\boxed{\theta = \arctan\left( \frac{\sum_{(x,y) \in \Omega} \sin\theta(x,y)}{\sum_{(x,y) \in \Omega} \cos\theta(x,y)} \right)} \tag{3.20}
\end{equation}

Спрощена форма (середнє арифметичне кутів):

\begin{equation}
\theta = \frac{1}{N_\Omega} \sum_{(x,y) \in \Omega} \arctan\left( \frac{v(x,y)}{u(x,y)} \right) \tag{3.21}
\end{equation}

\subsubsection{Модель просторового положення}

Координати центру об'єкта в площині зображення:

\begin{align}
x_c &= \frac{x_{\min} + x_{\max}}{2} \tag{3.22} \\
y_c &= \frac{y_{\min} + y_{\max}}{2} \tag{3.23}
\end{align}

Глибина (Z-координата) визначається як медіанне значення карти глибини в області об'єкта:

\begin{equation}
\boxed{z = \text{median}\left\{ D(x,y) \mid (x,y) \in \Omega \right\}} \tag{3.24}
\end{equation}

Альтернативно — як середнє арифметичне:

\begin{equation}
z = \frac{1}{|\Omega|} \iint_{\Omega} D(x,y) \, dx\,dy \tag{3.25}
\end{equation}

\subsubsection{Модель траєкторії об'єкта}

Траєкторія об'єкта представляється у вигляді дискретної множини положень у часі:

\begin{equation}
\boxed{T(t) = \left\{ p(t_1), p(t_2), \ldots, p(t_n) \right\}} \tag{3.26}
\end{equation}

де $p(t_i) = (x_i, y_i)$ — положення об'єкта в момент $t_i$.

Для зменшення шумів використовується згладжена траєкторія (ковзне середнє):

\begin{equation}
\tilde{p}(t_i) = \frac{1}{2k+1} \sum_{j=-k}^{k} p(t_{i+j}) \tag{3.27}
\end{equation}

\subsubsection{Модель прискорення об'єкта}

Прискорення об'єкта визначається як похідна швидкості за часом:

\begin{equation}
\boxed{a = \frac{v(t+1) - v(t)}{\Delta t}} \tag{3.28}
\end{equation}

де $\Delta t = 1$ кадр (дискретний час).

\subsection{Об'єднання моделей і методів ідентифікації}

\subsubsection{Сумісна модель GeoNet}

Модель GeoNet реалізує сумісну оцінку глибини та оптичного потоку на основі геометричних обмежень:

\begin{equation}
\boxed{I_2(p) = I_1(p + f(p))} \tag{3.29}
\end{equation}

\begin{equation}
f(p) = T(p) \cdot \frac{Z(p) - Z_0}{Z(p)} \cdot \text{proj}(p) \tag{3.30}
\end{equation}

Ця модель дозволяє використовувати синергію між двома типами даних.

\subsubsection{Ансамблева модель Bagging}

Ансамблева модель Bagging передбачає усереднення результатів кількох моделей:

\begin{equation}
\boxed{\hat{y}_{\text{bag}}(x) = \frac{1}{M} \sum_{i=1}^{M} f_i(x)} \tag{3.31}
\end{equation}

У контексті оцінки глибини:

\begin{equation}
D_{\text{bag}}(x,y) = \frac{1}{M} \sum_{i=1}^{M} D_i(x,y) \tag{3.32}
\end{equation}

\subsubsection{Адаптивна модель Boosting}

Модель Boosting використовує зважену комбінацію з виправленням помилок:

\begin{equation}
\hat{y}_{\text{boost}}(x) = \sum_{i=1}^{M} \alpha_i f_i(x) \tag{3.33}
\end{equation}

Вага кожної моделі обчислюється на основі її помилки:

\begin{equation}
\alpha_i = \frac{1}{2} \ln\left( \frac{1 - \epsilon_i}{\epsilon_i} \right) \tag{3.34}
\end{equation}

Адаптивне комбінування результатів оцінки глибини:

\begin{equation}
\boxed{D_{\text{final}}(x) = \begin{cases} 
D_1(x), & \text{якщо } |\nabla D_1(x)| \leq \tau \\ 
D_2(x), & \text{якщо } |\nabla D_1(x)| > \tau 
\end{cases}} \tag{3.35}
\end{equation}

\subsubsection{Інтегральна модель ідентифікації}

Інтегральна модель об'єднує всі компоненти в єдину систему:

\begin{equation}
\boxed{\Theta = \mathcal{P}\left(\mathcal{F}(I_1,I_2), \mathcal{Z}(I_2), \mathcal{D}(I_2)\right)} \tag{3.36}
\end{equation}

де:
\begin{itemize}
\item $\mathcal{F}$ — модель оптичного потоку;
\item $\mathcal{Z}$ — модель оцінки глибини;
\item $\mathcal{D}$ — модель виявлення об'єктів;
\item $\mathcal{P}$ — модель обчислення параметрів.
\end{itemize}

Повна система рівнянь для обчислення параметрів має вигляд:

\begin{equation}
\begin{cases}
\text{boxes} = \mathcal{D}(I_2) \\
F = \mathcal{F}(I_1, I_2) \\
D = \mathcal{Z}(I_2) \\
\forall \text{box} \in \text{boxes}: \quad v = \frac{1}{|\Omega|} \iint_{\Omega} \|F\| dxdy \\
\quad \theta = \arctan\left(\frac{\iint_{\Omega} \sin\theta(F) dxdy}{\iint_{\Omega} \cos\theta(F) dxdy}\right) \\
\quad z = \underset{(x,y)\in\Omega}{\text{median}} D(x,y)
\end{cases} \tag{3.37}
\end{equation}

\subsection{Висновки до розділу}

У розділі 3 розглянуто математичні моделі для дистанційної ідентифікації параметрів динамічних об'єктів. Проведений аналіз дозволяє зробити наступні висновки:

\begin{enumerate}
    \item \textbf{Концептуальна модель} (3.1) визначає загальну структуру системи ідентифікації як композицію модулів виявлення, оцінки оптичного потоку, оцінки глибини та обчислення параметрів (3.6).
    
    \item \textbf{Модель DETR} (3.7-3.11) забезпечує кінцево-кінцеве виявлення об'єктів з використанням механізму багатоголової уваги, що дозволяє моделювати довгострокові залежності в зображенні.
    
    \item \textbf{Моделі оцінки глибини MiDaS та DPT} (3.12-3.16) використовують багатомасштабні функції втрат та трансформерну архітектуру для отримання щільних карт глибини з високою точністю.
    
    \item \textbf{Моделі визначення параметрів} на основі оптичного потоку дозволяють обчислювати:
    \begin{itemize}
        \item швидкість об'єкта (3.17-3.18);
        \item напрямок руху (3.20-3.21);
        \item просторове положення (3.22-3.25);
        \item траєкторію (3.26-3.27);
        \item прискорення (3.28).
    \end{itemize}
    
    \item \textbf{Моделі об'єднання} (Bagging, Boosting, GeoNet) дозволяють підвищити точність та надійність ідентифікації шляхом комбінування різних методів. Особливо ефективною є адаптивна модель Boosting (3.35), яка використовує локальні характеристики зображення для вибору найкращого методу.
    
    \item \textbf{Інтегральна модель} (3.36-3.37) об'єднує всі компоненти в єдину систему, що забезпечує повний цикл дистанційної ідентифікації — від вхідного відеопотоку до вихідних параметрів об'єктів.
\end{enumerate}

Розроблені математичні моделі створюють теоретичну основу для практичної реалізації системи дистанційної ідентифікації параметрів динамічних об'єктів. У наступному розділі буде розглянуто програмну реалізацію цих моделей та експериментальні результати.

\end{document}